\section{Algorytmy lookahead} \label{sec:lookahead}

Algorytmy \textit{lookahead} stanowią prostszą alternatywę dla MCTS. Zamiast budować drzewo przez losowe symulacje, agent systematycznie sprawdza wszystkie możliwe sekwencje akcji do zadanej głębokości i ocenia końcowy stan za pomocą heurystycznej funkcji oceny~\cite{dockhorn2019}. Podejście to nie wymaga kosztownych \textit{rolloutów}, jednak liczba stanów do sprawdzenia rośnie wykładniczo z głębokością --- co w Hearthstone, z jego dużą przestrzenią akcji na turę, szybko staje się problemem.

\subsection{Greedy: 1-Step Lookahead}

Najbardziej podstawową formą jest agent \textbf{zachłanny}, który rozważa jedynie bezpośrednie konsekwencje każdej legalnej akcji i wybiera tę prowadzącą do najwyżej ocenionego stanu po jednym ruchu (ang.~\textit{1-step lookahead}). W terminologii SabberStone odpowiada to domyślnemu \textit{greedy agent} opisanemu w sekcji~\ref{sec:baseline}. Boty tego typu dominowały liczbowo w konkursie w latach 2018--2020, lecz rzadko zajmowały czołowe miejsca --- w edycji 2020 najwyższy klasyfikowany bot zachłanny zajął 7. miejsce (win-rate 0{,}570 w UCDT)~\cite{hearthstone_competition}.

\subsection{N-Step Lookahead}

Rozszerzeniem jest przeszukiwanie na głębokość $N$ ruchów. Agent symuluje wszystkie dostępne sekwencje akcji do głębokości $N$ i ocenia liście drzewa funkcją heurystyczną. W edycji 2019 pierwsze miejsce w ścieżce \textit{Premade Deck Playing Track} zajął bot \textbf{Heimbrodta} oparty na Dynamic Lookahead, a trzecie --- bot Öfelego i Mendeza Ruiza z Greedy Lookahead~\cite{hearthstone_competition}. W edycji 2020 algorytmy lookahead wyraźnie zdominowały konkurencję, zajmując 7 z 10 pierwszych miejsc:
\begin{itemize}
\item 1.\ miejsce: Sebastian Miller, \textit{2-Step Lookahead}, win-rate 0{,}723,
\item 2.\ miejsce: Bohnhof, Graeffe, \textit{Dynamic Lookahead}, win-rate 0{,}709,
\item 4.\ miejsce: Tracht, Mozheiko, \textit{3-Step Lookahead}, win-rate 0{,}676,
\item 5.\ miejsce: Markus Hempel, Julia Heise, \textit{2-Step Lookahead}, win-rate 0{,}660,
\item 6.\ miejsce: Kartik Mudgal, Kushagra Kumar, \textit{DFS Lookahead}, win-rate 0{,}586,
\item 7.\ miejsce: Kartik Mudgal, Kushagra Kumar, \textit{DFS Lookahead}, win-rate 0{,}578, 
\item 9.\ miejsce: Sathya Sudha Murugan, Maik Buettner, \textit{Balanced Lookahead}, win-rate 0{,}541~\cite{hearthstone_competition}.
\end{itemize}

\subsection{Dynamic Lookahead}

Wariantem adaptacyjnym jest \textbf{Dynamic Lookahead}, w którym głębokość przeszukiwania jest dostosowywana na bieżąco w zależności od dostępnego czasu obliczeniowego lub złożoności aktualnego stanu planszy. W ścieżce UCDT 2020 pierwsze i drugie miejsce zajęły boty Kiasa i Babela (win-rate 0{,}655) oraz Bohnhofa i Graeffego (win-rate 0{,}632), oba oparte na Dynamic Lookahead z talią MidrangeJadeShaman~\cite{hearthstone_competition}. Dynamiczne zarządzanie głębokością pozwala unikać przekroczenia limitu czasu na ruch, zachowując jednocześnie możliwie największą głębokość przeszukiwania.

\subsection{Lookahead z modelowaniem przeciwnika}

Uzupełnieniem algorytmów lookahead jest modelowanie ruchów przeciwnika, opisane przez Dockhorna i in.~\cite{dockhorn2018}. Zamiast oceniać jedynie własne akcje, agent symuluje równocześnie odpowiedź przeciwnika w kolejnej turze i ocenia stan gry po obu ruchach. Dzięki temu agent unika ruchów, które wyglądają korzystnie w izolacji, lecz pozostawiają łatwe do wykorzystania pozycje. Bot autorstwa Dockhorna wygrał ścieżkę \textit{Premade Deck Playing Track} w 2018 roku właśnie z zastosowaniem tej techniki~\cite{dockhorn2018, hearthstone_competition}.

\subsection{Lookahead a MCTS — porównanie}

Zarówno MCTS, jak i lookahead wymagają dostępu do modelu przejść (ang.~\textit{forward model}) silnika gry. Kluczowa różnica polega na tym, że MCTS buduje drzewo probabilistycznie, kierując zasoby obliczeniowe na obiecujące gałęzie, podczas gdy lookahead przeszukuje drzewo deterministycznie i wyczerpująco do zadanej głębokości. Wyniki konkursów z lat 2018--2020 pokazują, że w Hearthstone lookahead z głębokością 2--3 jest co najmniej równoważny MCTS przy porównywalnym budżecie obliczeniowym, a w edycji 2020 konsekwentnie go przewyższał~\cite{hearthstone_competition, dockhorn2019}.
